% Options for packages loaded elsewhere
\PassOptionsToPackage{unicode}{hyperref}
\PassOptionsToPackage{hyphens}{url}
\documentclass[
  11pt,
  a4paper,
]{article}
\usepackage{xcolor}
\usepackage[margin=1in]{geometry}
\usepackage{amsmath,amssymb}
\setcounter{secnumdepth}{-\maxdimen} % remove section numbering
\usepackage{iftex}
\ifPDFTeX
  \usepackage[T1]{fontenc}
  \usepackage[utf8]{inputenc}
  \usepackage{textcomp} % provide euro and other symbols
\else % if luatex or xetex
  \usepackage{unicode-math} % this also loads fontspec
  \defaultfontfeatures{Scale=MatchLowercase}
  \defaultfontfeatures[\rmfamily]{Ligatures=TeX,Scale=1}
\fi
\usepackage{lmodern}
\ifPDFTeX\else
  % xetex/luatex font selection
\fi
% Use upquote if available, for straight quotes in verbatim environments
\IfFileExists{upquote.sty}{\usepackage{upquote}}{}
\IfFileExists{microtype.sty}{% use microtype if available
  \usepackage[]{microtype}
  \UseMicrotypeSet[protrusion]{basicmath} % disable protrusion for tt fonts
}{}
\makeatletter
\@ifundefined{KOMAClassName}{% if non-KOMA class
  \IfFileExists{parskip.sty}{%
    \usepackage{parskip}
  }{% else
    \setlength{\parindent}{0pt}
    \setlength{\parskip}{6pt plus 2pt minus 1pt}}
}{% if KOMA class
  \KOMAoptions{parskip=half}}
\makeatother
\setlength{\emergencystretch}{3em} % prevent overfull lines
\providecommand{\tightlist}{%
  \setlength{\itemsep}{0pt}\setlength{\parskip}{0pt}}
\usepackage{bookmark}
\IfFileExists{xurl.sty}{\usepackage{xurl}}{} % add URL line breaks if available
\urlstyle{same}
\hypersetup{
  pdftitle={Knowledge Graph Report},
  hidelinks,
  pdfcreator={LaTeX via pandoc}}

\title{Knowledge Graph Report}
\author{}
\date{}

\begin{document}
\maketitle

\section{Knowledge Graph Report}\label{knowledge-graph-report}

\subsection{Introduction}\label{introduction}

This project develops a knowledge graph for London public transportation
with a particular focus on the Transport for London (TfL) network,
Oyster/contactless fare semantics, accessibility, operators, night
services, and journey-planning facts. The domain choice aligns directly
with the coursework brief\textquotesingle s recommended "public
transportation" theme and supports a mixture of operational,
accessibility, and fare-related competency questions.

The repository implements a hybrid construction workflow. A manually
curated ontology and seed ABox are stored in
\texttt{ontologies/manual/tfl\_kamyar\_final.ttl}, the structured
pipeline is implemented in
\texttt{data\_retrieval/structured\_retrieval/tfl\_api\_processor.py},
the unstructured pipeline is implemented across
\texttt{data\_retrieval/unstructured\_retrieval/}, and the merge step is
orchestrated by \texttt{main.py}. The final merged artifact in the
repository is
\texttt{ontologies/pipeline\_output/final\_london\_transport\_kg.ttl},
which currently contains 17,539 triples and occupies roughly 656 KB on
disk.

The public repository linked to this submission is:

\begin{itemize}
\tightlist
\item
  \href{https://github.com/cobogo124/sadiqkhanfanclub}{sadiqkhanfanclub}
\end{itemize}

The important outcome is that the graph covers the
coursework\textquotesingle s required question set while also exposing
where the current build is still incomplete. The current repository
snapshot is strong on documentation of intent, source manifests, and
competency-question coverage, but it is also a hybrid snapshot rather
than a fully clean rebuild from raw inputs. In particular, the manually
curated ontology currently provides most of the queryable
\texttt{http://example.org/tfl\#} facts used by the coursework SPARQL
queries, while the automated outputs require a namespace-alignment fix
before they can contribute seamlessly to the same query space.

\subsection{Data Source Selection}\label{data-source-selection}

\subsubsection{Structured sources}\label{structured-sources}

The primary structured source is the official TfL Unified API, cached
locally in \texttt{downloads/tfl\_api\_cache.json}. The committed cache
is approximately 79 MB and contains 695 line records across bus, Tube,
Overground, DLR, and Elizabeth line services, with 7,949 nested
stop-point entries. This choice is appropriate for the coursework
because it offers authoritative operational identifiers, line metadata,
route membership, and station/stop attributes in a machine-readable
form.

The repository also includes a generated source manifest in
\texttt{docs/generated/tfl-sources.json} and
\texttt{docs/generated/tfl-sources.md}. The JSON manifest records 1,299
retrievals across the following categories:

\begin{itemize}
\tightlist
\item
  \texttt{core}: 1
\item
  \texttt{textual}: 2
\item
  \texttt{route\_sequences}: 695
\item
  \texttt{line\_stop\_points}: 19
\item
  \texttt{line\_statuses}: 19
\item
  \texttt{search}: 3
\item
  \texttt{stop\_points}: 560
\end{itemize}

This is valuable from a reproducibility point of view because it shows
the intended breadth of the structured collection process, not just the
small subset ultimately materialized into the main queryable namespace.

As a secondary structured reference, the repository includes
\texttt{downloads/naptan\_london.csv} (about 4.3 MB). NaPTAN is useful
as a normalization and identifier-support resource for stops and
stations, particularly where the TfL API returns multiple stop-point
variants or where station naming ambiguity arises.

\subsubsection{Textual sources}\label{textual-sources}

The unstructured pipeline begins from a curated list of 21 Wikipedia
articles defined in
\texttt{data\_retrieval/unstructured\_retrieval/wiki\_scraper.py}. These
articles cover ticketing, accessibility, night services, network
context, and general transport semantics, including entries such as
"Oyster card", "London fare zones", "Step-free access", "Night Tube",
"Night buses in London", and "Interchange station".

The source manifest additionally records two non-Wikipedia textual
resources that are closer to official evidence:

\begin{itemize}
\tightlist
\item
  TfL step-free guide PDF
\item
  TfL Freedom of Information page for London Overground operator
  information
\end{itemize}

This hybrid source strategy is sensible for the domain. Official
structured APIs are the best fit for line, stop, and status data, while
textual sources are useful for extracting explanatory semantics such as
accessibility, concessions, and service concepts that do not always
appear in a clean schema from the API.

\subsubsection{Source-selection
rationale}\label{source-selection-rationale}

The data-source selection balances four considerations:

\begin{enumerate}
\def\labelenumi{\arabic{enumi}.}
\tightlist
\item
  Authority. TfL API responses and official TfL pages are the most
  defensible evidence for operational transport facts.
\item
  Coverage. Wikipedia articles broaden the semantic range to include
  accessibility, fare-policy, and service concepts.
\item
  Reproducibility. Cached JSON, a manifest of retrieval URLs, and
  committed Turtle outputs reduce dependence on live services during
  marking.
\item
  Coursework fit. The combination satisfies the requirement to use at
  least one structured source and one textual source.
\end{enumerate}

\subsubsection{Source-related caveats}\label{source-related-caveats}

Two limitations should be acknowledged explicitly.

First, the repository currently preserves the top-level API cache and
the generated source manifest, but not the full
\texttt{downloads/raw/...} snapshot tree referenced by the manifest.
Second, the unstructured pipeline refers to raw Wikipedia text snapshots
under \texttt{data/raw/...}, but those raw text files are not committed
in the current repository snapshot. The graph outputs and LLM cache are
present, but the full end-to-end provenance trail is therefore
incomplete.

\subsection{Extension Of Existing
Ontologies}\label{extension-of-existing-ontologies}

The core ontology in \texttt{ontologies/manual/tfl\_kamyar\_final.ttl}
imports and extends two existing ontologies:

\begin{itemize}
\tightlist
\item
  GTFS (\texttt{http://vocab.gtfs.org/terms\#})
\item
  Schema.org (\texttt{http://schema.org/})
\end{itemize}

This satisfies the coursework requirement to reuse existing ontologies
and extend them with domain-specific classes and properties. The main
benefits are interoperability and avoiding the unnecessary reinvention
of general transport and web-data concepts.

\subsubsection{Class extensions}\label{class-extensions}

The ontology introduces several subclasses rooted in GTFS and Schema.org
concepts. Important examples include:

\begin{itemize}
\tightlist
\item
  \texttt{tfl:TfLLine\ rdfs:subClassOf\ gtfs:Route}
\item
  \texttt{tfl:TfLStation\ rdfs:subClassOf\ gtfs:Station}
\item
  \texttt{tfl:OysterFareZone\ rdfs:subClassOf\ gtfs:Zone}
\item
  \texttt{tfl:TfLOperator\ rdfs:subClassOf\ gtfs:Agency}
\item
  \texttt{tfl:OysterFare\ rdfs:subClassOf\ schema:PriceSpecification}
\item
  \texttt{tfl:Journey\ rdfs:subClassOf\ schema:Trip}
\end{itemize}

These are then specialized further into subclasses such as
\texttt{UndergroundLine}, \texttt{NightTubeLine}, \texttt{BusRoute},
\texttt{InterchangeStation}, \texttt{ElizabethLineStation},
\texttt{PublicToilet}, and \texttt{CarPark}. This layered design is
appropriate because it keeps the ontology recognizably
transport-oriented while allowing the graph to answer the
coursework\textquotesingle s more domain-specific questions.

\subsubsection{Property extensions}\label{property-extensions}

The ontology also defines useful transport-specific properties and
aligns them upward to GTFS and Schema.org where appropriate. Key
examples include:

\begin{itemize}
\tightlist
\item
  \texttt{tfl:servedByLine\ rdfs:subPropertyOf\ gtfs:route}
\item
  \texttt{tfl:operatesInZone\ rdfs:subPropertyOf\ gtfs:zone}
\item
  \texttt{tfl:routeNumber\ rdfs:subPropertyOf\ gtfs:shortName}
\item
  \texttt{tfl:lineColour\ rdfs:subPropertyOf\ gtfs:color}
\item
  \texttt{tfl:operatedBy\ rdfs:subPropertyOf\ schema:provider}
\item
  \texttt{tfl:journeyOrigin\ rdfs:subPropertyOf\ schema:departureStation}
\item
  \texttt{tfl:journeyDestination\ rdfs:subPropertyOf\ schema:arrivalStation}
\item
  \texttt{tfl:fareAmount\ rdfs:subPropertyOf\ schema:price}
\end{itemize}

These choices are technically strong because they allow local modeling
decisions to remain compatible with broader vocabularies. They also
support clearer reporting and potential reuse outside the coursework.

\subsubsection{Domain-specific concepts}\label{domain-specific-concepts}

The ontology adds concepts that are important for this domain but not
cleanly provided by the reused ontologies, including:

\begin{itemize}
\tightlist
\item
  fare concessions
\item
  night services
\item
  station facilities
\item
  accessibility features
\item
  engineering closures
\item
  alternative services
\item
  multi-leg journeys
\end{itemize}

This is a good example of extending rather than replacing external
ontologies: GTFS and Schema.org provide a skeleton, while the TfL
ontology adds the domain semantics needed by the competency questions.

\subsection{Mappings}\label{mappings}

\subsubsection{Structured mapping
pipeline}\label{structured-mapping-pipeline}

The structured pipeline is implemented in
\texttt{data\_retrieval/structured\_retrieval/tfl\_api\_processor.py}.
Its workflow is:

\begin{enumerate}
\def\labelenumi{\arabic{enumi}.}
\tightlist
\item
  Load a local JSON cache if available, otherwise fetch live TfL API
  data.
\item
  Iterate through each line and attach stop-point data.
\item
  Map each line to a domain class according to \texttt{modeName}.
\item
  Create station resources from stop points.
\item
  Interpret \texttt{additionalProperties} as either station facilities
  or accessibility features.
\item
  Serialize the result to
  \texttt{ontologies/pipeline\_output/structured\_london\_transport.ttl}.
\end{enumerate}

The mode-to-class mapping is straightforward and domain-appropriate:

\begin{itemize}
\tightlist
\item
  \texttt{tube} -\textgreater{} \texttt{UndergroundLine}
\item
  \texttt{dlr} -\textgreater{} \texttt{DLRLine}
\item
  \texttt{elizabeth-line} -\textgreater{} \texttt{ElizabethLineService}
\item
  \texttt{overground} -\textgreater{} \texttt{OvergroundLine}
\item
  \texttt{bus} -\textgreater{} \texttt{BusRoute}
\end{itemize}

At station level, the pipeline maps stop-point metadata into ontology
facts such as:

\begin{itemize}
\tightlist
\item
  \texttt{servedByLine}
\item
  \texttt{hasFacility}
\item
  \texttt{hasAccessibilityFeature}
\item
  \texttt{hasStepFreeStreetToPlatform}
\end{itemize}

This is a reasonable mapping design because it translates
source-specific JSON fields into domain semantics that directly support
the competency questions.

\subsubsection{Unstructured mapping
pipeline}\label{unstructured-mapping-pipeline}

The unstructured workflow is distributed across three scripts:

\begin{itemize}
\tightlist
\item
  \texttt{wiki\_scraper.py}
\item
  \texttt{extract\_triples.py}
\item
  \texttt{triples\_to\_rdf.py}
\end{itemize}

The process is:

\begin{enumerate}
\def\labelenumi{\arabic{enumi}.}
\tightlist
\item
  Download text from selected Wikipedia pages.
\item
  Split text into chunks.
\item
  Run spaCy named-entity recognition to identify relevant candidate
  entities.
\item
  Use an LLM prompt, or cached LLM output, to extract candidate triples.
\item
  Filter the extracted triples against recognized entities.
\item
  Map free-text predicates to ontology properties.
\item
  Serialize the resulting RDF and attach provenance links.
\end{enumerate}

The repository currently ships with \texttt{USE\_LLM\ =\ False}, meaning
the build depends on cached triple extraction results stored in
\texttt{data/caches/triples\_cache.json}. That cache currently contains
563 entries. This is useful for reproducibility and offline marking,
even though it means the present build is not using a live LLM call at
runtime.

The predicate-mapping stage in \texttt{triples\_to\_rdf.py} is
particularly important. It normalizes a wide range of surface predicates
such as \texttt{operatedBy}, \texttt{hasStop}, \texttt{acceptedOn},
\texttt{hasAccessibilityFeature}, and \texttt{openedDate} into ontology
properties under the TfL namespace. The script also adds
\texttt{dcterms:source} and \texttt{prov:wasGeneratedBy}, which gives
the unstructured portion of the graph better provenance than many
student projects achieve.

\subsubsection{Merge step}\label{merge-step}

\texttt{main.py} loads the manual ontology, loads the GTFS ontology,
runs the unstructured pipeline, runs the structured pipeline, and
serializes the merged graph to
\texttt{ontologies/pipeline\_output/final\_london\_transport\_kg.ttl}.

In design terms, this is a good architecture for the coursework because
it separates:

\begin{itemize}
\tightlist
\item
  ontology definition
\item
  structured extraction
\item
  unstructured extraction
\item
  final graph assembly
\end{itemize}

\subsubsection{Important integration
limitation}\label{important-integration-limitation}

The current repository snapshot also exposes the main technical weakness
of the build. The automated outputs do not currently share the same
namespace as the manual ontology.

\begin{itemize}
\tightlist
\item
  The structured output serializes resources under a Linux file-based
  \texttt{file:///.../ontologies/pipeline\_output/tfl\#} namespace
\item
  The unstructured output serializes resources under a Windows
  file-based \texttt{file:///.../ontologies/pipeline\_output/tfl\#}
  namespace
\item
  The manual ontology and all coursework SPARQL queries use
  \texttt{http://example.org/tfl\#}
\end{itemize}

This means the repository currently contains three parallel TfL
namespaces rather than one canonical graph namespace. In practice, this
is why the coursework SPARQL queries are answered primarily by the
manually curated \texttt{http://example.org/tfl\#} ABox rather than by
the automated outputs. The automated pipeline still adds a large amount
of data, but it is not yet fully integrated into the namespace used by
the competency-question layer.

This limitation should be documented rather than hidden, because it is
the single clearest remaining blocker to claiming a fully automated
end-to-end KG population workflow.

\subsection{Queries}\label{queries}

The repository includes 20 competency questions and their SPARQL
implementations in
\texttt{competency\_questions/competency\_question.txt}. The executable
validation harness is
\texttt{competency\_questions/competency\_question\_test.py}, which
loads
\texttt{ontologies/pipeline\_output/final\_london\_transport\_kg.ttl}
and runs all 20 questions.

Running that validation script against the current repository snapshot
returns non-empty results for all 20 questions, so the recorded CQ
answerability score is currently:

\begin{itemize}
\tightlist
\item
  20/20 answered
\item
  100\% non-empty query coverage
\end{itemize}

Representative results from the current graph include:

\begin{itemize}
\tightlist
\item
  King\textquotesingle s Cross St Pancras is linked to Circle,
  Hammersmith \& City, Metropolitan, Northern, Piccadilly, and Victoria
  lines.
\item
  The Victoria line terminates at Brixton and Walthamstow Central.
\item
  Stratford is modeled as an interchange served by Central, Jubilee,
  DLR, Elizabeth line, and Overground services.
\item
  Baker Street has a public toilet facility.
\item
  The Night Tube lines returned by the graph are Central, Jubilee,
  Northern, Piccadilly, and Victoria.
\item
  Three fare concessions for disabled bus passengers are represented and
  queryable.
\end{itemize}

It is also worth noting that the repository contains two slightly
different CQ descriptions: one in \texttt{README.md} and one in
\texttt{competency\_questions/competency\_question.txt}. For the
purposes of implementation and validation, the authoritative set is the
latter, because it is the one paired with executable SPARQL and a
working test script.

\subsection{Evaluation Methodology}\label{evaluation-methodology}

\subsubsection{Quality metrics}\label{quality-metrics}

The most direct quality metric for this coursework is
competency-question answerability. On the current repository snapshot,
all 20 stored SPARQL queries return at least one answer.

Additional useful quality indicators from the current graph are:

\begin{itemize}
\tightlist
\item
  final merged graph size: 17,539 triples
\item
  queryable \texttt{http://example.org/tfl\#} station-like instances: 30
\item
  queryable line-like instances under the core ontology hierarchy: 39
\item
  queryable facility-like instances: 7
\item
  queryable accessibility-feature-like instances: 19
\item
  operators: 4
\item
  fare zones: 6
\item
  fare concessions: 3
\item
  subjects with explicit \texttt{dcterms:source}: 211
\item
  distinct Wikipedia source URIs represented in the graph: 15
\end{itemize}

These numbers indicate that the graph is not merely a minimal toy
example. It contains enough domain structure to answer fare,
accessibility, operator, closure, and journey questions across multiple
modes.

\subsubsection{Performance metrics}\label{performance-metrics}

A lightweight performance evaluation was run directly against the
committed final graph. On this environment:

\begin{itemize}
\tightlist
\item
  parsing \texttt{final\_london\_transport\_kg.ttl} took about 0.80
  seconds
\item
  running the full 20-query validation set took about 0.28 seconds in
  total
\item
  average per-query time was about 0.014 seconds
\item
  peak resident memory during the run was about 49 MB
\end{itemize}

The slowest query in this run was the King\textquotesingle s Cross
intersection query, but even that completed in roughly 0.12 seconds.
These timings are acceptable for a coursework-sized KG and suggest that
the current graph remains easy to load and query on ordinary hardware.

\subsubsection{Provenance and source-coverage
evaluation}\label{provenance-and-source-coverage-evaluation}

The generated source manifest is a useful secondary metric because it
shows how much source material the pipeline was designed to cover:

\begin{itemize}
\tightlist
\item
  1,299 manifest entries in \texttt{docs/generated/tfl-sources.json}
\item
  695 route-sequence fetches
\item
  560 stop-point fetches
\item
  19 line-status fetches
\end{itemize}

This provides a stronger argument for scalability than the queryable
manual ABox alone. However, the namespace split described earlier means
that much of this harvested information is not yet visible in the main
coursework query namespace.

\subsubsection{Suggested baseline
comparisons}\label{suggested-baseline-comparisons}

The coursework brief also asks for comparison against simpler LLM
approaches. The repository does not preserve a full baseline experiment,
so the most defensible methodology is:

\begin{enumerate}
\def\labelenumi{\arabic{enumi}.}
\tightlist
\item
  Use the 20 competency questions as a fixed evaluation set.
\item
  Record gold or expected answers from the current validated KG.
\item
  Ask a plain LLM the same 20 questions without retrieval.
\item
  Ask the same LLM with KG-backed retrieval.
\item
  Compare exactness, completeness, hallucination rate, and citation
  quality.
\end{enumerate}

This comparison is not fully archived in the current repository, so it
should be presented as the evaluation method rather than as a completed
experimental result.

\subsubsection{Validity threats}\label{validity-threats}

The evaluation should acknowledge four threats to validity:

\begin{enumerate}
\def\labelenumi{\arabic{enumi}.}
\tightlist
\item
  The automated outputs are currently split across file-based namespaces
  and are therefore underused by the main SPARQL layer.
\item
  The repository no longer contains the complete raw source snapshot
  tree referenced in the manifest.
\item
  The unstructured extraction stage relies on cached LLM output rather
  than a fully reproducible live run.
\item
  Some CQ answers are currently supported by manually seeded ABox facts
  in the ontology, which improves answerability but reduces the strength
  of the end-to-end automation claim.
\end{enumerate}

\subsection{Conclusion}\label{conclusion}

Overall, this repository contains a strong coursework narrative: it
defines an ontology for TfL public transport, reuses existing
vocabularies appropriately, combines structured and unstructured
sources, and validates the resulting KG with 20 SPARQL competency
questions that all return answers on the current final artifact.

The main remaining weakness is not the domain framing, but integration
quality. The project is closest to a successful submission when
described honestly as a hybrid system: a solid manually curated core
ontology and CQ layer, supported by an ambitious automated pipeline
whose harvested outputs still need namespace normalization and tighter
provenance preservation before they can be counted as fully integrated
automated population of the same graph.

\end{document}
